% ===================================================================
%  Αρχείο: 1737_SOLUTION.tex (Fragment)
%  Αυτόματη μετατροπή μέσω doc2latex.py
% ===================================================================

ΘΕΜΑ 4

Θεωρούμε ορθογώνιο τρίγωνο ΑΒΓ ($\widehat{Α}$ = 90\textsuperscript{ο}) και το ύψος του ΑΗ. Ονομάζουμε Δ και Ε τα συμμετρικά σημεία του Η ως προς τις ευθείες ΑΒ και ΑΓ αντίστοιχα. Αν Μ είναι το σημείο τομής του τμήματος ΗΔ με την πλευρά ΑΒ και Ν είναι το σημείο τομής του ΗΕ με την πλευρά ΑΓ, να αποδείξετε ότι:

\begin{alist}
\item ΑΗ=ΑΔ=ΑΕ. \monades{10}

\item Η γωνία ΕΗΔ είναι ορθή. \monades{8}

\item Τα σημεία Α, Ε και Δ είναι συνευθειακά και ΜΝ = $\frac{\Delta Ε}{2}$ . \monades{7}
\end{alist
